%% ---------------------------------------------------------------- Fig. 1
\begin{figure*}[!htbp]
\centering
\noindent\makebox[\linewidth][l]{\textbf{a}}\\[1mm]
\resizebox{\linewidth}{!}{%
\begin{tikzpicture}[
  font=\footnotesize,
  box/.style={draw, rounded corners=2pt, minimum height=9mm, minimum width=22mm,
              align=center, inner sep=2pt},
  arr/.style={-{Latex[length=2mm]}, thick},
  lbl/.style={font=\scriptsize, midway, above=0.5mm, inner sep=0pt}]
\node[box, fill=black!5] (dev) at (0,0) {devices};
\node[box, fill=edgec!15, draw=edgec] (edge) at (4.1,0)
     {edge tier\\[-1pt]\scriptsize $n_e\times$ class};
\node[box, fill=gwc!15, draw=gwc] (gw) at (8.2,0)
     {gateway tier\\[-1pt]\scriptsize $n_g\times$ class};
\node[box, fill=cloudc!12, draw=cloudc] (cl) at (12.3,0)
     {cloud\\[-1pt]\scriptsize class};
\draw[arr] (dev) -- node[lbl] {access link} (edge);
\draw[arr] (edge) -- node[lbl] {access link} (gw);
\draw[arr] (gw) -- node[lbl] {WAN} (cl);
\draw[arr, dashed] (dev.south) to[out=-35, in=-145]
      node[font=\scriptsize, midway, below=0.5mm] {alternative stage placement} (gw.south);
\node[font=\scriptsize, text=edgec, above=1mm of edge] {replication $r$};
\end{tikzpicture}}

\vspace{6mm}
\noindent\makebox[\linewidth][l]{\textbf{b}}\\[1mm]
\resizebox{\linewidth}{!}{%
\begin{tikzpicture}[
  font=\scriptsize,
  n/.style={circle, draw, minimum size=4mm, inner sep=0pt},
  pr/.style={rectangle, draw=red!70!black, fill=red!8, rounded corners=1pt,
             align=center, inner sep=2pt},
  ev/.style={rectangle, draw=edgec, fill=edgec!12, rounded corners=1pt,
             align=center, inner sep=2pt},
  e/.style={-, thin}, el/.style={font=\scriptsize, midway, sloped, above=-0.5mm}]
\node[n] (root) at (7,3.6) {};
\node[n] (a) at (3.2,2.4) {}; \node[n] (b) at (10.8,2.4) {};
\draw[e] (root) -- node[el] {$n_e=2$} (a);
\draw[e] (root) -- node[el] {$n_e=8$} (b);
\node[pr] (p1) at (0.9,1.0) {latency bound\\$>$ deadline};
\node[n]  (c)  at (5.2,1.0) {};
\draw[e] (a) -- node[el] {low-end class} (p1);
\draw[e] (a) -- node[el] {high-end class} (c);
\node[ev] (e1) at (3.9,-0.5) {evaluate};
\node[pr] (p2) at (7.0,-0.5) {bound dominated\\by evaluated design};
\draw[e] (c) -- node[el] {$r=1$} (e1);
\draw[e] (c) -- node[el] {$r=3$} (p2);
\node[pr] (p3) at (9.6,1.0) {cost bound\\$>$ budget};
\node[ev] (e2) at (12.7,1.0) {evaluate};
\draw[e] (b) -- node[el] {fast link} (p3);
\draw[e] (b) -- node[el] {standard link} (e2);
\node[font=\scriptsize, anchor=west] at (0.0,3.6)
     {\begin{tabular}{@{}l@{}}partial\\designs\end{tabular}};
\end{tikzpicture}}
\caption{\textbf{Exploration by elimination.} \textbf{a}, Architecture template: the
number and class of edge nodes and gateways, the class of the cloud instance, the access
link class, the replication factor of the ingest stage and the placement of the
processing stages are all design decisions. \textbf{b}, Each node of the traversal is a
partially specified architecture. It is discarded when an optimistic bound on a
constrained quantity already violates its limit (red), or when its optimistic objective
vector is dominated by an architecture that has already been evaluated (red); only
designs that survive to a leaf reach the expensive evaluator (blue).}
\label{fig:ce:concept}
\end{figure*}
